\noindent\textbf{ZUR HANNOVERSCHEN TRIANGULATION.}\hfill347
\begin{center}\textsc{Gauss} an \textsc{Schumacher}.\quad Göttingen, 10. September 1818.\end{center}

Ich eile Ihnen anzuzeigen, dass ich von unserm Minister \textsc{Arnswaldt} den
Auftrag erhalten, die zur Verbindung einer hannoverschen Triangulirung mit
der Ihrigen nöthigen Messungen in Lüneburg vorzunehmen und dazu das
Nöthige mit Ihnen zu verabreden. Er macht zugleich mir Hoffnung, dass
demnächst auch die Fortsetzung selbst wohl zu Stande kommen werde, und
es freut mich, dass diese nun durch die in Lüneburg vorzunehmenden Operationen
gesichert werden kann.\ .\ .\ .\ .\ .\ .

\begin{center}\textsc{Gauss} an \textsc{Schumacher}.\quad Göttingen, 20. Mai 1820.\end{center}
\begingroup
\makeatletter
\renewcommand{\@makefnmark}{\hbox{\normalfont[\textsuperscript{*})]}}
\long\def\@makefntext#1{\parindent1.5em\noindent[\textsuperscript{*})\ #1]}
\makeatother
\par\noindent\hspace*{\parindent}.\ .\ .\ .\ .\ .\ In dieser Ungewissheit\footnote{Über \textsc{Schumachers} Aufenthalt.} adressire ich diesen Brief nach Copenhagen
und wünsche sehnlich, dass er Sie treffen und \textls[120]{bald} treffen möge: er
soll Ihnen nemlich die Nachricht anzeigen, dass in Folge eines Schreibens
vom Grafen \textsc{von Münster} aus London, als Antwort meines vor einem Jahre
\makeatletter
\renewcommand{\@makefnmark}{\hbox{\normalfont[\textsuperscript{**})]}}
\long\def\@makefntext#1{\parindent1.5em\noindent[\textsuperscript{**})\ #1]}
\makeatother
von Ihnen gefälligst besorgten Briefes\footnote{Abgedruckt in Band IV, S.~482/483.},
\begin{quote}
»der König die Fortsetzung der Gradmessung durch das Königreich
Hannover genehmigt hat«.
\end{quote}
\par\medskip
.\ .\ .\ .\ .\ .
\endgroup

\begin{center}\textsc{Gauss} an \textsc{Schumacher}.\quad Göttingen, 4. März 1821.\end{center}

.\ .\ .\ .\ .\ .\ Es war früher meine Absicht, von Hamburg anzufangen und so
von Norden nach Süden zu messen, allein da ich leider so sehr durch die
Schreibfaulheit und Unzuverlässigkeit aller Künstler, mit denen ich zu thun
habe, hingehalten werde, und bis jetzt noch gar nichts von den nöthigen
Hülfsmitteln in Händen habe, so würden die daraus erwachsenden Verlegenheiten
noch viel grösser sein, wenn ich jenen Plan befolgte. Dieser und noch
\textls[120]{verschiedene andere wichtige} Gründe nöthigen und bewegen mich zu
dem umgekehrten Plan, von Süden nach Norden zu messen.\ .\ .\ .\ .\ .\ .
\par
\makeatletter
\insert\footins{\reset@font\footnotesize\hsize\columnwidth\parindent0pt
\smallskip\hbox to\hsize{\hfil44\textsuperscript{*}}\par}
\makeatother

\clearpage
\noindent348\hfill\textbf{BRIEFWECHSEL MIT SCHUMACHER.}
\begin{center}\textsc{Gauss} an \textsc{Schumacher}.\quad Wulfsode, 18. September 1822.\end{center}

.\ .\ .\ .\ .\ .\ So viel eine vorläufige Inspection des Terrains urtheilen lässt,
würde es nicht unmöglich sein, \textls[120]{die ganze Linie} von Breithorn bis Eschede
[Scharnhorst] (11220 Meter lang) unmittelbar zu messen. Welch eine herrliche
Basis wäre dies!\ .\ .\ .\ .\ .\ .

\begin{center}\textsc{Gauss} an \textsc{Schumacher}.\quad [Göttingen, Januar 1825.]\end{center}

.\ .\ .\ .\ .\ .\ Wenn ich alle grössern und kleinern Durchhaue aus den Jahren
1821---1824 zusammen zähle, von solchen, wo vielleicht ein Dutzend Bäume
gefällt sind, bis zu den grössten, so mögen etwa 16 oder 17 Durchhaue vorgekommen
sein. Der Allergrösste, nach der Ausdehnung, war im Becklinger
Holz unweit der Strasse von Bergen nach Soltau.\ .\ .\ .\ .\ .\ .

\begin{center}\textsc{Gauss} an \textsc{Schumacher}.\quad Dangast, 1 Stunde von Varel, 20. Junius 1825.\end{center}

.\ .\ .\ .\ .\ .\ Ich selbst pflege [beim Heliotrop] durch die drei Fussspitzen
einen Kreis zu beschreiben, dessen Centrum als Zielpunkt betrachtet wird.
Bei meinen beiden neuesten Heliotropen ist noch das Centrum selbst durch eine
Spitze bezeichnet, welches viel Bequemlichkeit verschafft.\ .\ .\ .\ .\ .\ .

\begin{center}\textsc{Gauss} an \textsc{Schumacher}.\quad Göttingen, 14. Januar 1829.\end{center}

Da wir noch wenig darüber wissen, ob die Figur der Erde von der Figur
des mittlern Ellipsoids in längern oder in kürzern Undulationen abweicht, so
bin ich jedenfalls der Meinung, dass man am besten thut, immer das mittlere
Ellipsoid zum Grunde zu legen.

Wenn ich jedoch es nicht gerade unbedingt verwerfen will, einmal ein
osculirendes Ellipsoid zum Grunde zu legen, so kann ich dies doch nur da
für zulässig halten, wo man Mittel hat, ein osculirendes Ellipsoid zu bestimmen,
d.\ i.\ ein solches, in dem die Krümmung \textls[120]{sowohl} im Sinne des Meridians,
als in dem darauf senkrechten Sinn der wirklichen Gestalt so nahe wie
möglich kommt. Dies ist aber bloss aus der Verbindung von \textls[120]{Messungen}
in beiderlei Sinn zu erhalten (Breiten- und Längen-Gradmessung).\ .\ .\ .\ .\ .\ .

\clearpage
\noindent\textbf{ZUR HANNOVERSCHEN TRIANGULATION.}\hfill349
\par\medskip
\begingroup
\makeatletter
\renewcommand{\@makefnmark}{\hbox{\normalfont\textsuperscript{*})}}
\makeatother
Wenn ein Meridian nicht wirklich elliptisch ist, so kann man allerdings
eine Ellipse berechnen, die sich an zwei Stücke eines Bogens anschliesst, und
man mag dies meinethalben eine osculirende Ellipse nennen; allein durch
Umdrehung dieser Ellipse um \textls[120]{ihre} kleine Axe\footnote{Von der man nur sagen kann, dass sie der Erdaxe parallel ist, ohne mit ihr zusammen zu fallen, ja von der sie in der Regel weit abstehen wird.} entsteht keine Fläche, die
man osculirendes Ellipsoid nennen darf, oder mit andern Worten, zwischen
dem wirklichen Werth des Längengrads und demjenigen Längengrad, den man
auf dem durch Umdrehung jener osculirenden Ellipse [entstandenen Ellipsoid]
berechnet, ist gar kein Zusammenhang. Durch Verwechselung beider setzt
man sich den grössten Fehlern aus.\ .\ .\ .\ .\ .\ .
\endgroup
\begin{center}\rule{.26\linewidth}{.4pt}\end{center}
